% We have proposed a novel method for instrumental variable regression, Deep Feature IV (DFIV), which performs two-stage least squares regression on flexible and expressive features of the instrument and treatment. As a contribution to the IV literature, we showed how to adaptively learn these feature maps with deep neural networks. We also showed that the off-policy policy evaluation (OPE) problem in deep RL can be interpreted as a nonlinear IV regression, and that DFIV performs competitively in this domain. In RL problems with additional external confounders, we believe DFIV will be of great value. This work brings the research worlds of deep offline RL and causality from observational data closer.

We have proposed a novel method for instrumental variable regression, Deep Feature IV (DFIV), which performs two-stage least squares regression on flexible and expressive features of the instrument and treatment. As a contribution to the IV literature, we showed how to adaptively learn these feature maps with deep neural networks. We also showed that the off-policy policy evaluation (OPE) problem in deep RL can be interpreted as a nonlinear IV regression, and that DFIV performs competitively in this domain. This work thus brings the research worlds of deep offline RL and causality from observational data closer.

In terms of future work, it would be interesting to adapt the ideas from \citep{angrist1995split,angrist1999jackknife,hansen2014instrumental} to select the regularization hyperparameters of DFIV as well as investigate generalizations of DFIV beyond the additive model \eqref{eq:structuralcausalmodel} as considered in \citep[Section 5.5]{Carrasco2007}. In RL, problems with additional confounders are common, see e.g. \citep{namkoong2020off,shang2019environment}, and we believe that adapting DFIV to this setting will be of great value.

%Furthermore, our work suggests that learning adaptive feature maps is a powerful tool not only in IV regression but also in general supervised regression problems. 
